% =============================================================================
% Hardened patch: §4 Finite-scale local blocks and whitening
%   sec:finite-scale-local-blocks-and-whitening
% =============================================================================

% =============================================================================
\section{Finite-scale local blocks and whitening}
\label{sec:finite-scale-local-blocks-and-whitening}
\statuspaper{LIVE}\quad
\statusmig{migrated}\quad
\statusreferee{in-progress}\quad
\statussympy[rh/sympy/sec-finite-scale-local-blocks-and-whitening.py]{verified}\quad
\statussim[rh/simulation/sec-finite-scale-local-blocks-and-whitening.py]{verified}\quad
\statuslean{n/a}
\par\medskip

The jet-limit identities of Section~\ref{sec:jet-limit-local-blocks}
produce the infinitesimal same-point and cross-point blocks.  The
comparison theorems of Section~\ref{sec:projective-rigidity} require the
corresponding two-center blocks at nonzero center separation, together
with a controlled normalization by the same-point blocks.  This section
records those finite-center blocks and the associated whitening loss.  No
source-energy transfer or actual-zeta suppression input is used here.

\begin{definition}[Finite packet domain]
\label{def:finite-packet-domain}
For a retained theta chart \(I_T\), define \(t_\pm=m\pm s/2\) and
\[
        \mathcal D_T
        :=\{(m,s)\in\mathbb R^2:
             t_-\in I_T,\ t_+\in I_T\}.
\]
Let \(\mathcal D_T^\times:=\{(m,s)\in\mathcal D_T:s\ne0\}\).  The
orientation convention is that rows correspond to the center \(t_-\) and
columns to the center \(t_+\).  In applications \(s>0\), but allowing
signed \(s\) makes the removable-singularity statement at \(s=0\) precise.
Uniform estimates below are taken on \(\mathcal D_T\), after the global
lower-height cutoff has been enlarged if necessary.
\end{definition}

\subsection{Finite same-point and mixed blocks}
\label{subsec:finite-same-point-and-mixed-blocks}

Fix \((m,s)\in\mathcal D_T^\times\).  Set
\[
        t_\pm=m\pm s/2,
        \qquad q_\pm=q(t_\pm),
        \qquad \Del_m(s)=\Ph(t_-)-\Ph(t_+).
\]

\begin{definition}[Finite same-point and mixed blocks]
\label{def:finite-same-point-and-mixed-blocks}
The finite same-point blocks are
\begin{equation}
\label{eq:Gpm}
G_{m,\pm}(s)
=\frac1\pi
\begin{pmatrix}
2q(t_\pm) & \tfrac12 q'(t_\pm)\\[0.8ex]
\tfrac12 q'(t_\pm) & \tfrac1{12}\bigl(q''(t_\pm)+2q(t_\pm)^3\bigr)
\end{pmatrix}.
\end{equation}
For \(s\ne0\), the finite mixed block is
\begin{equation}
\label{eq:Nm}
N_m(s)
=\frac1\pi
\begin{pmatrix}
-\dfrac{2\sin\Del_m(s)}{s}
&
\dfrac{\sin\Del_m(s)+q_+ s\cos\Del_m(s)}{s^2}
\\[2.4ex]
-\dfrac{\sin\Del_m(s)+q_- s\cos\Del_m(s)}{s^2}
&
\dfrac{(q_-+q_+)s\cos\Del_m(s)+(2-q_-q_+s^2)\sin\Del_m(s)}{2s^3}
\end{pmatrix}.
\end{equation}
\end{definition}

\begin{lemma}[Exact finite-block interpretation]
\label{lem:finite-blocks-exact-interpretation}
For \((m,s)\in\mathcal D_T^\times\),
\[
        G_{m,\pm}(s)=J(t_\pm).
\]
Moreover,
\begin{equation}
\label{eq:Nm-kernel-derivative-form}
N_m(s)=
\begin{pmatrix}
2K_\Ph(t_-,t_+) & \partial_y K_\Ph(t_-,t_+)\\[0.6ex]
\partial_x K_\Ph(t_-,t_+) & \tfrac12\partial_{xy}K_\Ph(t_-,t_+)
\end{pmatrix}.
\end{equation}
Equivalently, if \(N_{12}\) is the matrix of \eqref{eq:cross-N12} with
\((T_1,T_2)=(t_-,t_+)\), then \(N_m(s)=\pi^{-1}N_{12}\), with
\(s_{12}=T_1-T_2=-s\) absorbed into the signs in \eqref{eq:Nm}.
\end{lemma}

\begin{proof}
The identity \(G_{m,\pm}(s)=J(t_\pm)\) is immediate from
\eqref{eq:same-point-J}.  For the mixed block, apply the off-diagonal
kernel derivative formulas of Lemma~\ref{lem:phase-kernel-derivatives}
at \((T_1,T_2)=(t_-,t_+)\).  There
\(s_{12}=t_- - t_+=-s\), \(q_1=q_-\), \(q_2=q_+\), and
\(\Delta=\Del_m(s)\).  Substitution gives the four entries of
\eqref{eq:Nm-kernel-derivative-form}, which are exactly the entries in
\eqref{eq:Nm}.
\end{proof}

\begin{lemma}[Removable singularity of \(N_m(s)\)]
\label{lem:Nm-removable-singularity}
If \(\Ph\in C^3\) on a neighborhood of \(m\), then the entries of
\(N_m(s)\) in \eqref{eq:Nm} extend continuously to \(s=0\).  The limiting
value is
\begin{equation}
\label{eq:Nm-at-zero}
        N_m(0)=J(m)
        = \frac1\pi
        \begin{pmatrix}
        2q(m) & \tfrac12 q'(m)\\[0.8ex]
        \tfrac12 q'(m) & \tfrac1{12}\bigl(q''(m)+2q(m)^3\bigr)
        \end{pmatrix}.
\end{equation}
\end{lemma}

\begin{proof}
Write
\[
        q=q(m),\qquad a=q'(m),\qquad b=q''(m).
\]
Taylor expansion about \(m\), using the symmetric points
\(m\pm s/2\), gives
\[
        \Del_m(s)
        =\Ph(m-s/2)-\Ph(m+s/2)
        =-qs-\frac{b}{24}s^3+o(s^3),
\]
\[
        q_+=q+\frac a2s+\frac b8s^2+o(s^2),
        \qquad
        q_-=q-\frac a2s+\frac b8s^2+o(s^2).
\]
Consequently
\[
        \sin\Del_m(s)
        =-qs+\left(\frac{q^3}{6}-\frac b{24}\right)s^3+o(s^3),
        \qquad
        \cos\Del_m(s)=1-\frac{q^2}{2}s^2+o(s^2).
\]
Substituting these expansions into \eqref{eq:Nm} gives
\begin{align*}
        -\frac{2\sin\Del_m(s)}{s}
        &=2q+o(1),\\
        \frac{\sin\Del_m(s)+q_+s\cos\Del_m(s)}{s^2}
        &=\frac a2+o(1),\\
        -\frac{\sin\Del_m(s)+q_-s\cos\Del_m(s)}{s^2}
        &=\frac a2+o(1),\\
        \frac{(q_-+q_+)s\cos\Del_m(s)+(2-q_-q_+s^2)\sin\Del_m(s)}{2s^3}
        &=\frac b{12}+\frac{q^3}{6}+o(1).
\end{align*}
Multiplication by the common factor \(1/\pi\) gives
\eqref{eq:Nm-at-zero}.
\end{proof}

\begin{remark}[Separate coalescence computation]
\label{rem:Nm-coalescence-not-cross-limit}
Lemma~\ref{lem:cross-block-jet-limit} is a fixed-center-separation limit
in the sampling parameter \(h\).  Lemma~\ref{lem:Nm-removable-singularity}
is instead a coalescence statement in the center-separation parameter
\(s\).  The latter is not obtained by using
Lemma~\ref{lem:cross-block-jet-limit} uniformly as \(s\to0\); it is the
separate Taylor computation above.
\end{remark}

\subsection{Whitened finite block}
\label{subsec:whitened-finite-block}

\begin{definition}[Whitened finite block]
\label{def:whitened-finite-block}
Whenever \(G_{m,-}(s)\) and \(G_{m,+}(s)\) are positive definite, let
\(G_{m,\pm}(s)^{-1/2}\) denote the positive definite inverse square root
and set
\begin{equation}
\label{eq:omega-finite}
\widehat\Om_\z(s;m)
:= G_{m,-}(s)^{-1/2}\,N_m(s)\,G_{m,+}(s)^{-1/2}.
\end{equation}
\end{definition}

\begin{corollary}[Coalescence: \(\widehat\Om_\z(0;m)=I_2\)]
\label{cor:omega-coalescence}
For every retained midpoint \(m\) at sufficiently large height,
\(\widehat\Om_\z(s;m)\) extends continuously to \(s=0\), and
\[
        \widehat\Om_\z(0;m)=I_2.
\]
\end{corollary}

\begin{proof}
By Lemma~\ref{lem:same-point-gram-positivity}, \(J(m)\succ0\) for every
sufficiently large retained midpoint \(m\).  Hence, by continuity,
\(G_{m,\pm}(s)=J(m\pm s/2)\) remain positive definite for all sufficiently
small \(|s|\) with \((m,s)\in\mathcal D_T\).  The map
\(G\mapsto G^{-1/2}\) is continuous on the positive definite cone.
Lemma~\ref{lem:Nm-removable-singularity} gives \(N_m(0)=J(m)\), while
\(G_{m,\pm}(0)=J(m)\).  Therefore
\[
\widehat\Om_\z(0;m)
=J(m)^{-1/2}J(m)J(m)^{-1/2}=I_2.
\]
\end{proof}

\subsection{Whitening loss}
\label{subsec:whitening-loss}

\begin{lemma}[Operator-norm bound for \(G_{m,\pm}^{-1/2}\)]
\label{lem:whitening-loss}
For \(\Ph_T=\theta\) on \(I_T\), after increasing the global lower-height
cutoff if necessary, \(G_{m,\pm}(s)\) are positive definite for every
\((m,s)\in\mathcal D_T\).  Moreover there is a quantity
\(\varepsilon_T\to0\), independent of \((m,s)\in\mathcal D_T\), such that
\begin{equation}
\label{eq:whitening-loss-bound}
        \bigl\|G_{m,\pm}(s)^{-1/2}\bigr\|_{\op}
        \le
        \sqrt{\frac{\pi}{2q(t_\pm)}}\,(1+\varepsilon_T)
        \qquad (T\to\infty).
\end{equation}
In particular, after possibly increasing the same lower-height cutoff,
\begin{equation}
\label{eq:whitening-loss-polynomial}
        \bigl\|G_{m,\pm}(s)^{-1/2}\bigr\|_{\op}
        \le Q(T)^{C_W}
\end{equation}
uniformly on \(\mathcal D_T\), with the absolute exponent \(C_W=0\)
whenever \(Q(T)\ge1\).
\end{lemma}

\begin{proof}
By Lemma~\ref{lem:finite-blocks-exact-interpretation},
\(G_{m,\pm}(s)=J(t_\pm)\).  Since \((m,s)\in\mathcal D_T\) implies
\(t_\pm\in I_T\subset[T/2,2T]\), the theta derivative estimates of
Lemma~\ref{lem:theta-derivative-asymptotics} are uniform at \(t=t_\pm\).
Repeating the proof of Lemma~\ref{lem:same-point-gram-positivity} with
\(T\) replaced by an arbitrary \(t\in I_T\) gives, uniformly on \(I_T\),
\[
        \lambda_{\min}(J(t))
        =\frac{2q(t)}{\pi}\,(1+o(1)).
\]
Equivalently, for some \(\delta_T\to0\) independent of \(t\in I_T\),
\[
        \lambda_{\min}(J(t))
        \ge \frac{2q(t)}{\pi}(1-\delta_T).
\]
For large \(T\), \(1-\delta_T>0\), so the finite same-point blocks are
positive definite and
\[
        \bigl\|G_{m,\pm}(s)^{-1/2}\bigr\|_{\op}
        =\lambda_{\min}(G_{m,\pm}(s))^{-1/2}
        \le
        \sqrt{\frac{\pi}{2q(t_\pm)}}\,(1-\delta_T)^{-1/2}.
\]
Absorbing \((1-\delta_T)^{-1/2}\) into \(1+\varepsilon_T\) proves
\eqref{eq:whitening-loss-bound}.  Finally, the uniform asymptotic above together with
Lemma~\ref{lem:phase-derivative-lower-bound} gives
\(\lambda_{\min}(J(t))\to\infty\) uniformly for \(t\in I_T\).  After
increasing the lower-height cutoff we therefore have
\(\lambda_{\min}(J(t))\ge1\) uniformly on \(I_T\).  Thus
\(\|G_{m,\pm}(s)^{-1/2}\|_{\op}\le1\le Q(T)^0\) whenever \(Q(T)\ge1\),
which is \eqref{eq:whitening-loss-polynomial} with \(C_W=0\).
\end{proof}

\begin{remark}[Scope of Lemma~\ref{lem:whitening-loss}]
\label{rem:whitening-loss-scope}
Lemma~\ref{lem:whitening-loss} controls only the real-axis same-point
whitening factors attached to the theta phase.  It does not prove
actual-zeta transverse suppression, does not supply source-energy input,
and does not give complex-disk or holomorphic whitening estimates.  Any
stronger whitening statement needed for a later complex-analytic argument
must be introduced as a separate input or proved in a separate section.
\end{remark}
